\section{\ref{sec:synthesis}장. 기계 전체}

요약 장이므로 새 실험을 소개하지 않는다. 각 행은 그 명제를 다룬 장과 같은 출처에 기댄다. \nt{GLUT} 버스와 \nt{GABA}를 통한 흐름 제어는 확고하게 확립되어 있고, 브로드캐스트 채널의 역할과 그 공동 동작은 대부분 가설이다.

{\footnotesize
\setlength{\tabcolsep}{3pt}\renewcommand{\arraystretch}{1.15}
\begin{longtable}{>{\raggedright}p{0.5cm}>{\raggedright}p{4.0cm}>{\raggedright}p{5.6cm}>{\raggedright}p{2.3cm}>{\raggedright\arraybackslash}p{1.7cm}}
\toprule
No. & 명제 & 실험과 측정한 것 & 출처 & 상태 \\
\midrule\endfirsthead
\toprule
No. & 명제 & 실험과 측정한 것 & 출처 & 상태 \\
\midrule\endhead
1 & 뉴런 $8.6\cdot10^{10}$개에 제어 신호는 네 종류뿐이다~\eqref{eq:machine} & 뇌 균질액에서 세포핵 계수(등방성 분획법): 성인 남성의 뉴런은 $86.1\pm8.1$\,$\times10^9$개 & \cite{Azevedo2009} & 측정됨 \\
2 & 명제~\ref{prop:architecture}의 앞 두 층: 데이터는 \nt{GLUT} 주소 버스를 따라 흐르고, 연결의 부호는 송신자가 정하며, 흐름은 \nt{GABA} 뉴런이 이끌고 정규화한다(\ref{sec:neurontypes}, \ref{sec:gaba}장) & 뉴런당 주된 신경전달물질은 하나(측정 총설); 피드포워드 억제는 피라미드 뉴런이 입력을 합산하는 창을 좁힌다; 전체 활동에 의한 나눗셈은 여러 영역에서 측정됨 & \cite{Strata1999,Pouille2001,Carandini2012} & 측정됨 \\
3 & 소자는 신뢰할 수 없다: 시냅스는 스파이크 대부분을 잃는다(표~\ref{tab:computer}, \ref{sec:code}장) & 해마 절편, 최소 자극: 스파이크의 절반 넘게 CA1에서 응답을 일으키지 못한다; 문턱값과 축삭 전도의 실패는 배제되었고, 원인은 확률적 신경전달물질 방출 & \cite{AllenStevens1994} & 측정됨 \\
4 & 뇌는 $\sim20$와트로 동작하며, 뉴런당 평균 초당 약 한 번 스파이크한다(\ref{sec:code}장); 엔지니어는 아직 이런 기계를 만들지 못했다 & 측정된 비용에서 얻은 추정~\eqref{eq:energy}: 시냅스 작업을 포함해 스파이크당 ATP 분자 약 $10^9$개(설치류 피질); 피질에 대한 상세 추정은 더 낮은 발화율을 준다. 기술 현황은 총설에 따름 & \cite{Attwell2001,Lennie2003,MehonicKenyon2022} & 이론 \\
5 & 대부분 시냅스의 학습은 국소적 적격 흔적과 하나의 공통 수 $\delta$의 곱이다(\eqref{eq:machine}의 둘째 식, \ref{sec:dofa}장) & 원숭이의 \nt{DOPA} 뉴런은 보상 예측 오차에 응답한다; 선조체 절편에서 도파민은 글루탐산 입력 $0.3$--$2$\,s 뒤에 올 때만 가시를 키운다. 적격 흔적이 측정됨 & \cite{Schultz1997,Yagishita2014} & 측정됨 \\
6 & 출력마다 따로 가는 오차 전선은 운동 학습 회로에만 있다(\ref{sec:motorlearning}장) & 소뇌: 등반섬유 신호가 입력과 동시에 오면 그 입력이 약해진다; 원숭이에서 푸르키네 세포의 복합 스파이크는 운동 오차의 방향을 나르고 보정을 일으킨다 & \cite{Ito2001,Herzfeld2018} & 측정됨 \\
7 & 각 브로드캐스트 채널은 몇 개 선로의 다발이다(명제~\ref{prop:bundle}) & \nt{DOPA}: 같은 과제에서 서로 다른 뉴런이 보상, 운동, 자극 특징을 나른다; 빠른 오차와 느린 도파민 수준은 서로 다르다. \nt{SERO}, \nt{NORA}, \nt{ACET}에서는 이런 분해가 덜 연구됨 & \cite{Engelhard2019,Hamid2016,Mohebi2019} & 측정됨 \\
8 & \nt{SERO}: 수준 조절기이자, 가설상 지평 $\tau_\gamma$(\ref{sec:serocomputer}장) & 자기 억제: 자기 5-HT$_{1A}$ 수용체의 작용제가 솔기핵 세로토닌 뉴런을 억제한다. 측정됨. 지평은 이론에서 제안됨; 가장 강한 실험은 생쥐에서 이 뉴런의 광유전학적 활성화가 지연 보상의 기다림을 늘린다는 것 & \cite{Sprouse1987,Doya2002,Miyazaki2014} & 가설 \\
9 & \nt{NORA}: 이득 $g$가 탐색과 결정 사이에서 고른다(\ref{sec:nora}장) & 신호 대 잡음비 증가: 깨어 있는 원숭이의 청각 피질에서 노르아드레날린은 동종의 울음소리에 대한 응답보다 배경 활동을 더 강하게 억제한다. 측정됨; 곱셈적 이득은 모델; 탐색과 결정 사이의 선택에서의 역할은 이론 & \cite{Foote1975NE,ServanSchreiber1990,AstonJones2005} & 가설 \\
10 & \nt{ACET}: 쓰기와 읽기를 전환한다(\ref{sec:acet}장) & 세 작용(내부 연결 약화, 입력 강화, 가소성 촉진)은 측정됨; 모드 전환은 사람에서 스코폴라민 실험으로 간접 지지됨 & \cite{HasselmoBower1992,Gil1997,Atri2004} & 가설 \\
11 & 식~\eqref{eq:machine}은 기계 전체를 기술한다 & \ref{sec:glutcomputer}--\ref{sec:acet}장 모델의 요약이며, 각 부분은 자기 장에 기댄다; 전체로서는 실험으로 검증되지 않음 & 장 안의 유도 & 가설 \\
12 & 노브들은 공통 제어 없이 협조해 동작한다: 오차, 놀람, 쓰기, 복귀(그림~\ref{fig:timeline}) & 실험 없음: 정성적 도식. 개별 고리는 \nt{NORA}와 \nt{ACET}의 분업 이론, 그리고 사람에서 동공과 학습 속도의 관계 & \cite{YuDayan2005,Nassar2012} & 가설 \\
13 & 하나의 공통 신호에 의한 학습은 결과에 영향을 주는 뉴런이 많을수록 느리다(명제~\ref{prop:broadcastcost}) & 선형 네트워크 학습 곡선 계산: 출력 섭동에 의한 학습은 경사 하강법보다 출력 수만큼 느리다 & \cite{Werfel2005} & 이론 \\
14 & 피질이 다층 회로를 어떻게 조정하는지는 알려져 있지 않다; 후보는 스파이크 버스트, 뉴런 가지 안의 계산, 예측에 의한 자기 지도 학습 & 오차 운반자로서의 버스트는 모델; 피질 뉴런 상세 모델의 응답은 $5$--$8$층 네트워크만 재현했다. 모델; 사람 피질 뉴런 가지에서 XOR을 내는 칼슘 스파이크는 절편에서 측정됨; 자기 지도 학습은 증거 총설 & \cite{Payeur2021,Beniaguev2021,Gidon2020,Keller2018} & 가설 \\
\bottomrule
\end{longtable}}
